\begin{definition}[Local semi-convexity and semi-concavity]
Let $N\in\mathbb N$, let $\Omega\subset\mathbb R^N$ and let $w:\mathbb R^N\to\mathbb R$.
We say that $w$ is \emph{locally semi-convex on $\Omega$}, written
$\mathrm{LocallySemiConvexOn}\ \Omega\ w$, if for every $x\in\Omega$ there are $\rho>0$ and
$\Lambda\ge0$ such that the Euclidean ball $B_\rho(x)$ is contained in $\Omega$ and
\[
 y\ \longmapsto\ w(y)+\tfrac{\Lambda}{2}\,|y|^{2}
\]
is convex on $B_\rho(x)$. A function $w$ is \emph{locally semi-concave on $\Omega$} when
$-w$ is locally semi-convex on $\Omega$; no separate name is introduced for that notion,
and the phrase is always used in that sense below.

This is Definition Definition~2.5 of the paper with the semi-convexity constant allowed to
depend on the point. The localization is forced by the paper's own use of the notion and
is not a weakening that loses content where the paper uses a global constant. Property (4)
of Section Section~2.3 concludes only that $\psi\circ w$ and $w\circ\Theta$ are
semi-convex \emph{on compact subsets}, with constants depending on the local Lipschitz
constant of $w$, so no single $\Lambda$ serves on all of $\Omega$ for the composed
functions; and those composed functions are exactly the ones fed to Proposition
Proposition~3.3 in Theorem Theorem~1.3(iii) ($u\circ R_h$) and inside Proposition~3.3 itself
($\psi_\lambda\circ u$). Conversely Proposition~3.3 never uses more than a constant valid on a
ball around one point (Step 4, near $x_1$), so the local notion carries the whole
argument. A function that is semi-convex on $\Omega$ in the sense of Definition~2.5, with a
global $\Lambda$, is locally semi-convex on $\Omega$ when $\Omega$ is open.
\end{definition}
